\chapter{Partial Orders, Lattices and Fixed Points}\label{ch:posets}

Order theory is where a degree turns from computing to structure, and it is the mathematics
that Part~V uses to make the engine terminate. This chapter relearns it on finite examples ---
divisors of $12$, subsets of a three-element set --- and ends with the theorem about fixed points
of monotone maps that underlies every ``iterate until it stops'' in computing. The engine has a
small world for all of it, in which every command is a decision procedure and every theorem says
the decision means the definition.

\section{Partial orders}

\begin{definition}
  A partial order on a set $P$ is a relation $\le$ that is reflexive ($x \le x$), antisymmetric
  ($x \le y$ and $y \le x$ imply $x = y$) and transitive ($x \le y$ and $y \le z$ imply $x \le z$).
  Write $x < y$ for $x \le y$ and $x \neq y$.
\end{definition}

\begin{example}
  Divisibility on the divisors of $12$: $\{1, 2, 3, 4, 6, 12\}$ with $a \le b$ iff $a \mid b$.
  Inclusion on the subsets of $\{a, b, c\}$. The chain $0 < 1 < \cdots < n - 1$.
\end{example}

A relation is usually given by a few generating pairs --- $a < b$, $a < c$ --- and the order is
the reflexive-transitive closure: the smallest reflexive, transitive relation containing them.
On a finite set the closure is reached by repeatedly adding $(x, z)$ whenever $(x, y)$ and $(y, z)$
are present, and $|P|$ rounds suffice because a chain of intermediate points has at most $|P|$
elements. Antisymmetry is not automatic and must be checked; the generating pairs $a < b$,
$b < a$ close to a relation that is not a partial order.

\section{Hasse diagrams}

\begin{definition}[Covers]
  $y$ covers $x$, written $x \lessdot y$, if $x < y$ and there is no $z$ with $x < z < y$.
\end{definition}

\begin{theorem}
  In a finite poset, $x \le y$ iff there is a chain $x = z_0 \lessdot z_1 \lessdot \cdots \lessdot
  z_k = y$ (with $k = 0$ allowed).
\end{theorem}
\begin{proof}
  A chain of covers gives $x \le y$ by transitivity. Conversely, if $x < y$, choose a maximal
  chain $x = z_0 < z_1 < \cdots < z_k = y$ (one exists, since chains are bounded in length by
  $|P|$); each step is a cover, since anything strictly between two consecutive elements could be
  inserted.
\end{proof}

So the covers determine the order, and the Hasse diagram --- elements as points, covers as edges,
larger elements drawn higher --- is a complete picture of a finite poset. That is what the
notebook draws.

\section{Bounds, joins and meets}

\begin{definition}
  An upper bound of $S \subseteq P$ is a $u$ with $s \le u$ for all $s \in S$. The \emph{join}
  (least upper bound, supremum) $\bigvee S$ is an upper bound that is $\le$ every upper bound.
  Dually, lower bounds and the \emph{meet} $\bigwedge S$.
\end{definition}

\begin{theorem}
  A join, if it exists, is unique.
\end{theorem}
\begin{proof}
  If $u, u'$ are both least upper bounds then $u \le u'$ and $u' \le u$, so $u = u'$ by antisymmetry.
\end{proof}

\begin{definition}
  A lattice is a poset in which every pair has a join and a meet. A bounded lattice has a top
  $\top$ and a bottom $\bot$.
\end{definition}

Divisors of $12$ form a lattice: the join is the lcm, the meet the gcd. Subsets form a lattice
with union and intersection. The poset with generators $a < b$, $a < c$ is not a lattice: $b$ and
$c$ have no upper bound at all.

\section{Monotone maps and fixed points}

\begin{definition}
  $f : P \to P$ is monotone if $x \le y$ implies $f(x) \le f(y)$. A fixed point of $f$ is an $x$
  with $f(x) = x$.
\end{definition}

\begin{theorem}[Knaster--Tarski, finite form]
  Let $P$ be a finite lattice with bottom $\bot$ and $f$ monotone. Then $f$ has a least fixed
  point, and it is $f^n(\bot)$ for the first $n$ with $f^{n+1}(\bot) = f^n(\bot)$.
\end{theorem}
\begin{proof}
  The sequence $\bot \le f(\bot) \le f^2(\bot) \le \cdots$ is increasing: $\bot \le f(\bot)$
  since $\bot$ is below everything, and if $f^k(\bot) \le f^{k+1}(\bot)$ then applying $f$ gives
  $f^{k+1}(\bot) \le f^{k+2}(\bot)$ by monotonicity. An increasing sequence in a finite poset
  stabilizes within $|P|$ steps (it cannot strictly increase more than $|P| - 1$ times), and
  where it stabilizes, $f^{n+1}(\bot) = f^n(\bot)$ is a fixed point.

  It is the least: let $y$ be any fixed point. Then $\bot \le y$, and if $f^k(\bot) \le y$ then
  $f^{k+1}(\bot) \le f(y) = y$. So every term of the sequence is $\le y$, in particular
  $f^n(\bot)$.
\end{proof}

The sequence is the Kleene chain, and the theorem is the reason every ``repeat until nothing
changes'' computation --- reachability, dataflow analysis, the closure of the first section ---
gives the least answer. The general Knaster--Tarski theorem, for complete lattices, gives the
least fixed point as $\bigwedge \{x : f(x) \le x\}$ without iteration; the finite form is the
one that computes.

\section{The engine's world}

\begin{minted}{lean}
def closure (elems : List String) (gen : List (String × String)) : List (String × String)
def checkPartialOrder (P : Poset) : Option String   -- a witness of failure, or none
def covers (P : Poset) (x y : String) : Bool :=
  P.lt x y && !P.elems.any fun z => P.lt x z && P.lt z y
def upperBounds (P : Poset) (xs : List String) : List String :=
  P.elems.filter fun u => xs.all fun x => P.rel x u
def sup (P : Poset) (xs : List String) : Option String :=
  let ubs := upperBounds P xs
  ubs.find? fun u => ubs.all fun v => P.rel u v
def latticeFailure (P : Poset) : Option (String × String × String)
def monotoneFailure (P : Poset) (f : PMap) : Option (String × String) :=
  P.le.find? fun (x, y) => !P.rel (f.apply x) (f.apply y)
def iterate (P : Poset) (f : PMap) (start : String) : List String
\end{minted}

Each is a decision over lists, and each is the definition transcribed: \lc{covers} is the
definition of a cover, \lc{sup} searches the upper bounds for one below all the others,
\lc{monotoneFailure} looks for a violating pair. The theorems say that the decisions mean the
definitions.

\begin{minted}{lean}
theorem checkPartialOrder_none {P : Poset} (h : checkPartialOrder P = none) :
    (∀ x ∈ P.elems, P.rel x x = true) ∧
    (∀ p ∈ P.le, p.1 ≠ p.2 → P.rel p.2 p.1 = false) ∧
    (∀ p ∈ P.le, ∀ q ∈ P.le, p.2 = q.1 → P.rel p.1 q.2 = true)

theorem covers_spec (P : Poset) (x y : String) :
    covers P x y = true ↔ P.lt x y = true ∧ ∀ z ∈ P.elems, ¬ (P.lt x z = true ∧ P.lt z y = true)

theorem sup_spec {P : Poset} {xs : List String} {j : String} (h : sup P xs = some j) :
    j ∈ upperBounds P xs ∧ ∀ v ∈ upperBounds P xs, P.rel j v = true
theorem sup_none {P : Poset} {xs : List String} (h : sup P xs = none) :
    ∀ j ∈ upperBounds P xs, ∃ v ∈ upperBounds P xs, P.rel j v = false

theorem monotone_of_none {P : Poset} {f : PMap} (h : monotoneFailure P f = none) :
    ∀ p ∈ P.le, P.rel (f.apply p.1) (f.apply p.2) = true

theorem iter_le_fixed {P : Poset} {f : PMap} (hm : monotoneFailure P f = none) {y : String}
    (hy : f.apply y = y) : ∀ (n : Nat) (x : String), P.rel x y = true → P.rel (f.iter n x) y = true
\end{minted}

The last is the second paragraph of the Knaster--Tarski proof, verbatim: by induction on $n$,
if $f^n(x) \le y$ then $f^{n+1}(x) \le f(y) = y$. With $x = \bot$ it says the chain stays below
every fixed point, so its last element --- checked to be a fixed point --- is the least. What the
engine does \emph{not} prove is the first paragraph, that the chain stabilizes within $|P|$
steps; \lc{iterate} runs at most $|P|$ rounds and the notebook checks that the last element is
fixed. That is the pigeonhole argument, and it is the chapter's second exercise.

\begin{logbox}[The shape of this world]
  Every theorem here has the form ``if the decision procedure says yes, the textbook Prop
  holds''. It is the shape of the run-time boundaries of Part~III, but here the Prop is the
  definition a reader meets in a first course. That is what should sit beside a Hasse diagram:
  not ``the engine drew this'', but ``these edges are exactly the covers'', with the theorem's
  name.
\end{logbox}

\begin{trybox}
\begin{minted}{text}
let D = divisors(12)
join(D, 4, 6)
meet(D, 4, 6)
lattice(D)
let N = poset({a,b,c}; a<b, a<c)
lattice(N)
let f = map(D; 1->2, 2->2, 3->6, 4->4, 6->6, 12->12)
monotone(D, f)
lfp(D, f)
le(D, 2, 12)
\end{minted}
  The joins in $D$ are least common multiples. $N$ fails with the witness pair. The least
  fixed point of $f$ is $2$, and the derivation is the Kleene chain $1 \to 2 \to 2$. The last
  line explains $2 \le 12$ by a chain of covers.
\end{trybox}

\section{Exercises}

\begin{exercisebox}[Closure]
  Prove that $|P|$ rounds of adding $(x, z)$ for $(x, y), (y, z)$ present compute the
  transitive closure on a set of size $|P|$, and find an example needing $\lceil \log_2 |P|
  \rceil$ rounds if each round composes the relation with itself instead.
\end{exercisebox}

\begin{exercisebox}[Pigeonhole, formally]
  State and prove in Lean that if $x_0 \le x_1 \le \cdots$ is a sequence in a finite poset with
  $x_{k+1} = f(x_k)$, then some $x_n = x_{n+1}$ with $n < |P|$. This closes the engine's one
  boundary in this world.
\end{exercisebox}

\begin{exercisebox}[Not a lattice]
  Draw the Hasse diagram of the subsets of $\{a, b, c\}$ with the empty set removed. Is it a
  lattice? Which pairs fail, and does \lc{latticeFailure} report the first of them in the
  engine's element order?
\end{exercisebox}

\begin{exercisebox}[Greatest fixed point]
  Dualize the theorem: starting from $\top$, the chain decreases to the greatest fixed point.
  Check that the engine's \texttt{gfp} is the dual code, and find a map with distinct least and
  greatest fixed points on the divisors of $12$.
\end{exercisebox}
